%% P23 --- Probability and Bayes' theorem
%%
%% Every number in this program that the reader cannot do in their head comes
%% from code/p23_probability_bayes.py and is pulled in with \val{}. The console
%% listing is a committed transcript written by that script. EVERYTHING in that
%% script is exact, over Fractions, and that is a decision rather than a
%% flourish: every claim this program makes is about the EQUALITY of two
%% expressions, and a tolerance would weaken each one into a resemblance.
%%
%% THE THESIS: a probability is a weight on a list of outcomes obeying three
%% rules; conditioning is changing what is underneath the fraction; and Bayes'
%% theorem is one identity written two ways. It is DERIVED here and never
%% asserted, because the derived form is the one you can rebuild at a desk and
%% the memorised form is the one people write upside down.
%%
%% WHAT THIS PROGRAM IS OWED, read out of the files rather than remembered:
%%   F10  hands this program over BY NAME, and its closing frame says exactly
%%        what the reader arrives with: "the numerator and the denominator, and
%%        the knowledge that CHOOSING THE DENOMINATOR WAS A DECISION". That is
%%        the on-ramp, because conditioning is that decision made deliberately.
%%        F10 also owns the union rule, which arrives here as the addition rule
%%        -- and the script gates this program against F10's own four committed
%%        counts, so the two cannot come apart about it.
%%   F13  owns the density, the fact that a density's HEIGHT is not a
%%        probability while its AREA is, and the shape of a weighted average.
%%   P12  owns the counting formally, so a naive probability's numerator and
%%        denominator are objects the reader can already compute.
%%
%% WHAT THIS PROGRAM LEAVES ALONE, checked against tools/programs.json:
%%   random variables, expectation and variance                    -> P24
%%   sums, the central limit theorem, concentration                -> P25
%%   estimation and maximum likelihood                             -> P26
%%   what a p-value is                                             -> P27
%%   KL and cross-entropy                                          -> P30
%%
%% Twin: programs/pl/P23-probability-bayes.tex. Frame for frame, macro for
%% macro, number for number; tools/parity.py compares them.

\program{Probability and Bayes' theorem}
\label{prog:P23}
\index{probability!as a measure}
\index{Bayes' theorem}

\begin{outcomes}
\outcome{1--8}{say what a sample space, an event and a probability are, and
  state the three rules a probability obeys}
\outcome{9--15}{write down what conditioning does to a space, and say what it
  needs in order to be defined}
\outcome{16--21}{derive Bayes' theorem from the definition rather than quoting
  it, and expand its denominator}
\outcome{22--30}{do the base-rate calculation, and say which of two
  conditionals a stated accuracy is}
\outcome{31--37}{say what independence means, and why it is not a property two
  events keep}
\outcome{38--47}{say what conditional independence assumes, and what a second
  piece of evidence is worth when it holds and when it does not}
\end{outcomes}

\begin{quiz}
\item Two events overlap. Is $\Prob(A \cup B) = \Prob(A) + \Prob(B)$?
  \teachesat{6--7}
  \answerto{No. Adding them counts the overlap twice, so
  $\Prob(A \cup B) = \Prob(A) + \Prob(B) - \Prob(A \cap B)$. The sum rule
  needs the events to be \emph{disjoint}, and that condition is the whole of
  what the third rule says.}
\item Write $\Prob(A \mid B)$ in terms of quantities that do not use the bar,
  and say what it needs in order to mean anything. \teachesat{10--15}
  \answerto{$\Prob(A \mid B) = \Prob(A \cap B) / \Prob(B)$, and it needs
  $\Prob(B) > 0$. It is a definition rather than a result: nothing is being
  proved, a new denominator is being chosen.}
\item Bayes' theorem falls out of one quantity written two ways. Which
  quantity? \teachesat{16--17}
  \answerto{The joint, $\Prob(A \cap B)$. The definition of conditioning gives
  it as $\Prob(A \mid B)\Prob(B)$ and also as $\Prob(B \mid A)\Prob(A)$;
  setting those equal and dividing is the whole derivation.}
\item You know $\Prob(B \mid A)$, $\Prob(B \mid \bar{A})$ and $\Prob(A)$. How
  do you get $\Prob(B)$? \teachesat{20--21}
  \answerto{Split $B$ by whether $A$ happened, which is a disjoint split, and
  add: $\Prob(B) = \Prob(B \mid A)\Prob(A) + \Prob(B \mid \bar{A})
  \Prob(\bar{A})$. That is the law of total probability, and it is the third
  rule applied to a split you chose.}
\item In odds form, what do you multiply the prior odds by to get the
  posterior odds? \teachesatone{41}
  \answerto{The likelihood ratio of the evidence:
  $\Prob(\text{evidence} \mid A) / \Prob(\text{evidence} \mid \bar{A})$.
  Written that way Bayes' theorem is a multiplication, which is why several
  pieces of evidence can be accumulated \dash{} provided the ratios are
  allowed to multiply, which is a condition and not a formality.}
\end{quiz}

% ============================================================
\section{A weight on a list of outcomes}
\label{sec:P23-measure}
\index{probability!sample space}
\index{probability!event}

\begin{fr}
Program~\ref{prog:F10} ends by handing this program a job. It gives you a
numerator, a denominator and one uncomfortable observation \dash{} that
\emph{choosing} the denominator was a decision somebody made \dash{} and then
says that what a probability actually \emph{is} belongs here.

So start where that program stopped. A request arrives at a service, is either
faulty or clean, and a fault detector either raises an alarm about it or stays
quiet. How many different things can happen to one request?
\dotline
\nextframe
\end{fr}

\begin{fr}
\ans{Four}

Faulty and flagged, faulty and missed, clean and flagged, clean and quiet.
Write them out, because the list is the object:

\begin{center}
\begin{tabular}{lll}
\toprule
& \textbf{alarm} & \textbf{quiet} \\
\midrule
\textbf{faulty} & caught & missed \\
\textbf{clean}  & false alarm & correctly ignored \\
\bottomrule
\end{tabular}
\end{center}

That list is the \textbf{sample space}. Two properties make it one: the
outcomes cover everything that can happen, and \emph{exactly} one of them
happens. Nothing else about it matters yet.

An \textbf{event} is a set of outcomes \dash{} which is Program~\ref{prog:F10}'s
object, arriving with a new name. Which outcomes make up the event
\emph{the detector raised an alarm}?
\dotline
\nextframe
\end{fr}

\begin{fr}
\ans{Two of them: \emph{caught} and \emph{false alarm}}

Not one. \emph{An alarm} is not the same event as \emph{a fault}, and the whole
of this program lives in the gap between those two sets.

A \textbf{probability} is a rule that hands each event a number. It obeys three
things and nothing else:

\begin{enumerate}
\item every event gets a number that is not negative;
\item the whole space gets $1$;
\item if two events cannot both happen, the number for \emph{either} is the sum
  of their two numbers.
\end{enumerate}

That is the entire definition. Everything in this program is a consequence of
those three, and it is worth noticing how little they say: nothing about
frequency, nothing about repetition, nothing about belief.

\mermaidfig{p23-space-and-measure}{The three objects: a list, a weighting of
it, and a set of its entries.}{outcome, event, weight}
\end{fr}

\begin{fr}
Take the first consequence now, because you can get it in one line. If $A$ is
an event, what is the probability that $A$ does \emph{not} happen?
\dotline
\nextframe
\end{fr}

\begin{fr}
\ans{$\Prob(\bar{A}) = 1 - \Prob(A)$}

$A$ and $\bar{A}$ cannot both happen, so rule three lets you add them; between
them they are the whole space, so rule two says the total is $1$. Two rules,
one line, and no picture was involved.

Keep that shape in mind. Every result in this program is obtained the same way
\dash{} split something into pieces that cannot overlap, add them, and use the
fact that the whole is one.
\end{fr}

\begin{fr}
Now the rule that gets broken. Program~\ref{prog:F10} measured two evaluation
sets: $\val{f10.eval.a}$ cases in one, $\val{f10.eval.b}$ in the other, and
$\val{f10.eval.shared}$ appearing in both, for $\val{f10.eval.union}$ distinct
cases altogether.

Draw one of those $\val{f10.eval.union}$ cases at random, so that every case is
equally likely. Is the probability that it is in the first set \emph{or} the
second equal to $\Prob(\text{first}) + \Prob(\text{second})$?
\dotline
\nextframe
\end{fr}

\begin{fr}
\begin{ansblock}
No. The two probabilities are $\frac{1}{2}$ and $\frac{3}{5}$, which add to
$\frac{11}{10}$, and no probability is above $1$.
\end{ansblock}

Rule three had a condition on it and the condition was doing work: the events
have to be unable to happen together. These two can, in
$\val{f10.eval.shared}$ cases out of $\val{f10.eval.union}$, so:
\[ \Prob(A \cup B) = \Prob(A) + \Prob(B) - \Prob(A \cap B) \]
and here that is $\frac{1}{2} + \frac{3}{5} - \frac{1}{10} = 1$, which is right
because every case is in one set or the other.

That is Program~\ref{prog:F10}'s counting rule with every count divided by
$\val{f10.eval.union}$. It is not an analogy: the script for this program reads
that program's four committed counts and checks both statements against each
other, so the two cannot come apart.
\end{fr}

\begin{fr}
It is worth saying plainly what has and has not happened in this section.

You now have a \textbf{measure}: something that hands each set a size, never
negative, additive on pieces that do not overlap, with the whole thing coming
to a fixed total. Program~\ref{prog:F13} built one of those out of area, and
the constraint it put on a density \dash{} that the area under it is exactly
$1$ \dash{} is rule two, written for a continuous space.


Program~\ref{prog:F13} spends a section on the fact that a density's
\emph{height} can exceed $1$ while a probability cannot. Which of the three
rules is the one that does not apply to a height?
\dotline
\nextframe
\end{fr}

% ============================================================
\section{Conditioning is changing the denominator}
\label{sec:P23-conditioning}
\index{probability!conditional}
\index{probability!denominator}

\begin{fr}
\begin{ansblock}
The second. It bounds the measure of a \emph{set}, and a height is not the
measure of a set.
\end{ansblock}

A height is a rate \dash{} probability per unit of $x$ \dash{} and nothing in
the three rules constrains it. Narrow the interval and the height climbs
without limit while the area it encloses stays exactly $1$.

Now to work. Here are the four outcomes with weights on them. A fault happens once in
$\val{p23.prev.den}$ requests; the detector catches $\val{p23.sens.pct}$ per
cent of faults, and is quiet about $\val{p23.sens.pct}$ per cent of clean
requests. Out of $\val{p23.alarm.den.long}$ requests:

\begin{center}
\begin{tabular}{lrr}
\toprule
& \textbf{alarm} & \textbf{quiet} \\
\midrule
\textbf{faulty} & $\val{p23.alarm.real}$ & $1$ \\
\textbf{clean}  & $\val{p23.alarm.false}$ & $\val{p23.quiet.clean}$ \\
\bottomrule
\end{tabular}
\end{center}

Nothing in this section needs those numbers to be realistic. What it needs is
that they add to $\val{p23.alarm.den.long}$, which they do, so this is a sample
space with a measure on it.

Now somebody tells you the alarm fired. What has changed about the space?
\dotline
\nextframe
\end{fr}

\begin{fr}
\begin{ansblock}
Two of the four outcomes are gone. Only the \emph{alarm} column is still
available, and its two entries no longer add to the whole.
\end{ansblock}

That is what conditioning is, and it is worth resisting the temptation to make
it sound like more. You have been handed a smaller space, so you renormalise:
divide by the weight of what is left.

\[ \Prob(A \mid B) \;=\; \frac{\Prob(A \cap B)}{\Prob(B)} \]

The numerator keeps only the part of $A$ that survives, and the denominator is
the new whole.
\end{fr}

\begin{fr}
This is Program~\ref{prog:F10}'s observation made deliberate. That program says
a naive probability is two counts and a division, and that choosing what goes
underneath is a decision. Conditioning is that decision taken on purpose:
\emph{given} what I have been told, what is underneath is no longer everything.

One question about that formula decides how you should treat it for the rest of
the program. Is it a definition, or is it something that was proved?
\dotline
\nextframe
\end{fr}

\begin{fr}
\ans{A definition}

\begin{warning}
Nothing above is being proved and there is no derivation behind it to go and
find \dash{} somebody chose to give that quantity a name and a bar. It matters
two sections from now, because Bayes' theorem \emph{is} a theorem, and knowing
which of the two you are looking at tells you which one you are entitled to
rebuild from scratch when you cannot remember it.
\end{warning}

Try it on the table. Of the $\val{p23.alarm.den.long}$ requests, how many
produce an alarm, and what fraction is that?
\dotline
\nextframe
\end{fr}

\begin{fr}
\begin{ansblock}
$\val{p23.alarm.real} + \val{p23.alarm.false}$ of them, which is
$\frac{\val{p23.alarm.num}}{\val{p23.alarm.den}}$, or
$\val{p23.alarm.pct}$ per cent.
\end{ansblock}

An alarm is a rare event, which is what you would want. Hold on to the two
pieces it came from rather than the total: $\val{p23.alarm.real}$ of those
alarms are about a real fault and $\val{p23.alarm.false}$ are not.

\begin{note}
Notice that neither number was given to you. The detector's specification
describes what it does to a faulty request and what it does to a clean one, and
these two counts are what happens when you also say how many of each there
were. Section 4 is entirely about the consequences of that sentence.
\end{note}
\end{fr}

\begin{fr}
The definition gets used more often rearranged than as it stands, and the
rearrangement is worth doing yourself, because it is the line the next section
is built out of.

Multiply both sides of the definition by $\Prob(B)$. What do you get?
\dotline
\nextframe
\end{fr}

\begin{fr}
\ans{$\Prob(A \cap B) = \Prob(A \mid B)\,\Prob(B)$}

Which reads: to get both, get $B$, and then get $A$ out of what is left. That
is the \textbf{multiplication rule}, and it is the definition with the
denominator moved across.

One condition travels with all of this and it is the only one in the program.
$\Prob(B)$ must be greater than zero \dash{} you cannot condition on something
that cannot happen, because the space you would be renormalising against has no
weight in it.

\begin{trapbox}
\textbf{$\Prob(A \mid B)$ and $\Prob(B \mid A)$ are different numbers, and
usually very different ones.} They have different denominators, which is the
whole of what the bar controls.

The temptation is real because English does not distinguish them well.
\enquote{The probability of an alarm, given a fault} and \enquote{the
probability of a fault, given an alarm} differ by three words, and a spoken
sentence often drops them. Section 4 measures how far apart the two can get on
an entirely ordinary detector.
\end{trapbox}

\mermaidfig{p23-two-denominators}{What conditioning changes, and what it keeps.}{conditioning restricts}
\end{fr}

% ============================================================
\section{Bayes, in one line}
\label{sec:P23-bayes}
\index{Bayes' theorem!derived}

\begin{fr}
Everything needed for the theorem is now on the page, and the derivation is
short enough to do rather than read.

Using the multiplication rule from the last section, write the joint
probability $\Prob(A \cap B)$ in \emph{two} different ways.
\dotline
\nextframe
\end{fr}

\begin{fr}
\begin{ansblock}
\[ \Prob(A \cap B) = \Prob(A \mid B)\,\Prob(B)
   = \Prob(B \mid A)\,\Prob(A) \]
\end{ansblock}

The set $A \cap B$ does not know which way round you decomposed it, so the two
expressions are the same number. Divide both sides by $\Prob(B)$:

\[ \Prob(A \mid B) \;=\; \frac{\Prob(B \mid A)\,\Prob(A)}{\Prob(B)} \]

and that is Bayes' theorem. There is nothing else in it. One symmetric quantity
written two ways, and one division.
\end{fr}

\begin{fr}
It is worth being blunt about how ordinary that is, because the theorem has a
reputation the derivation does not support. Nothing was assumed, no limit was
taken, and no new object was introduced. What it does is \textbf{reverse a
conditional}: it takes the direction you can measure and gives you the
direction you wanted.

That is exactly the shape of most measurement problems in this field. A test
set tells you how the model behaves on cases you have labelled; what you wanted
was how a case behaves given the model's output.
\end{fr}

\begin{fr}
The four parts carry names, and the names are the least useful thing about
them. $\Prob(A)$ is the \textbf{prior}, $\Prob(B \mid A)$ the
\textbf{likelihood}, $\Prob(B)$ the \textbf{evidence} and $\Prob(A \mid B)$ the
\textbf{posterior}.

\begin{note}
Two of those names are worth a caution. The \emph{likelihood} is a conditional
probability read as a function of what is being conditioned on, which is not
what the English word suggests. And \emph{evidence} names the denominator,
which is the piece nobody computes directly \dash{} the next frame says what
you do instead. Neither name will help you rebuild the formula; the derivation
you have just done will.
\end{note}

Three of those four are usually handed to you by whoever specified the thing
you are reasoning about. Which one is not?
\dotline
\nextframe
\end{fr}

\begin{fr}
\ans{The evidence, $\Prob(B)$}

Nobody publishes it, and it is the part of the theorem people get stuck on. You
will usually know the two conditionals pointing the other way instead, and that
turns out to be enough.

You know $\Prob(B \mid A)$, $\Prob(B \mid \bar{A})$ and $\Prob(A)$. Write
$\Prob(B)$ in terms of them.
\dotline
\nextframe
\end{fr}

\begin{fr}
\begin{ansblock}
\[ \Prob(B) = \Prob(B \mid A)\,\Prob(A)
            + \Prob(B \mid \bar{A})\,\Prob(\bar{A}) \]
\end{ansblock}

Split $B$ according to whether $A$ happened. Those two pieces cannot overlap
and between them they are all of $B$, so rule three adds them \dash{} and each
piece is the multiplication rule again. This is the \textbf{law of total
probability}, and it is the third rule applied to a split you chose.

\textbf{It is also where the trouble in the next section hides.} That formula
contains $\Prob(A)$, so the answer depends on how common $A$ is, whether or not
anybody mentioned it.
\end{fr}

% ============================================================
\section{What fraction of the alarms are real}
\label{sec:P23-baserate}
\index{probability!base rate}
\index{troubleshooting!prosecutor's fallacy}

\begin{fr}
This section is the reason the program exists, and it opens with a question you
should answer at speed. Do not compute anything, and do not look back at the
table \dash{} the point is what you say before you work it out.

A fault detector is $\val{p23.sens.pct}$ per cent accurate, in the ordinary
sense: it catches $\val{p23.sens.pct}$ per cent of faults and is quiet about
$\val{p23.sens.pct}$ per cent of clean requests. Faults happen once in
$\val{p23.prev.den}$ requests. It has just raised an alarm.

Roughly what fraction of its alarms are real faults? Write down a number.
\dotline
\nextframe
\end{fr}

\begin{fr}
\begin{ansblock}
$\frac{\val{p23.ppv.num}}{\val{p23.ppv.den}}$, which is
$\val{p23.ppv.pct}$ per cent. Nine in ten of its alarms are false.
\end{ansblock}

\begin{trapbox}
\textbf{\enquote{It is $\val{p23.sens.pct}$ per cent accurate, so an alarm is
almost certainly a real fault.}} If you wrote down something near
$\val{p23.sens.pct}$, you are in the majority, and the reasoning that got you
there is not stupid \dash{} it is one substitution away from correct.

$\val{p23.sens.pct}$ per cent is $\Prob(\text{alarm} \mid \text{fault})$: the
detector's behaviour on the faulty requests. What you were asked for is
$\Prob(\text{fault} \mid \text{alarm})$: the composition of its alarms. The
last section's warning box named those as different numbers; this is how
different.

The reason is the denominator, and section 3's law of total probability put it
there in writing. Out of $\val{p23.alarm.den.long}$ requests there are
$\val{p23.alarm.real}$ real alarms and $\val{p23.alarm.false}$ false ones
\dash{} \textbf{about $\val{p23.alarm.ratio}$ false alarms for every real
one} \dash{} because a $\val{p23.fpr.pct}$ per cent error rate on an enormous
population of clean requests produces more alarms than a
$\val{p23.sens.pct}$ per cent catch rate on a tiny population of faulty ones.
\end{trapbox}

Now do it the other way, because the two routes have to agree and seeing that
they do is what makes the formula usable rather than decorative. Write the
answer using Bayes' theorem, with the denominator expanded by the law of total
probability.
\dotline
\nextframe
\end{fr}

\begin{fr}
\begin{ansblock}
\[ \Prob(F \mid A)
   = \frac{\Prob(A \mid F)\,\Prob(F)}
          {\Prob(A \mid F)\,\Prob(F) + \Prob(A \mid \bar{F})\,\Prob(\bar{F})}
   = \frac{\val{p23.ppv.num}}{\val{p23.ppv.den}} \]
\end{ansblock}

Every quantity in that fraction was in the detector's specification except one.
$\Prob(F)$, the rate at which faults happen at all, appears twice and nobody
mentioned it when the detector was described.

\begin{aibox}
\textbf{Precision} \dash{} the fraction of flagged cases that turn out to be
real \dash{} is $\Prob(\text{real} \mid \text{flagged})$, which is exactly the
quantity this section is about. \textbf{Recall} is
$\Prob(\text{flagged} \mid \text{real})$, which is the other one.

A classification report prints both, measured on a population somebody chose,
and that population is usually balanced, because a balanced set makes every
class visible. Production traffic is not. Recall conditions on the real class
alone, so it is a property of the model and it travels; precision conditions on
what the model flagged, so it depends on how many of each class there were and
it does not. Two adjacent lines of the same report, behaving completely
differently, with nothing on the card to say so.
\end{aibox}

\mermaidfig{p23-two-populations}{Where the alarms come from, and why counting
both matters.}{where alarms come from}
\end{fr}

\begin{fr}
The fallacy has a name, and it comes from a courtroom rather than from
engineering. A prosecutor tells a jury that the probability of the evidence
matching, given innocence, is one in a million, and invites them to conclude
that the probability of innocence, given the match, is one in a million.

Those are the two conditionals from section 2's warning box. Write down both of
them for this detector, as percentages: the fraction of \emph{clean requests}
that raise an alarm, and the fraction of \emph{alarms} that are about clean
requests.
\dotline
\nextframe
\end{fr}

\begin{fr}
\begin{ansblock}
$\Prob(\text{alarm} \mid \text{clean}) = \val{p23.fpr.pct}$ per cent and
$\Prob(\text{clean} \mid \text{alarm}) = \val{p23.false.pct}$ per cent.
\end{ansblock}

The detector almost never troubles a clean request, and almost every alarm it
raises is about one. Both are true at once, of the same detector, on the same
day \dash{} the first is a property of the detector and the second is a
property of the traffic.

Now vary the one quantity nobody stated. Same detector throughout \dash{}
nothing about it changes \dash{} and only the rate at which faults occur moves.

Holding the instrument fixed is what makes the next four rows an argument
rather than an illustration. Whatever the answers turn out to be, they cannot
be a property of the detector, because the detector is the one thing in the
table that did not change.

Before you read on, fill in the middle row. At a fault rate of one in a
hundred, a detector wrong one time in a hundred: what fraction of the alarms
are real?

\begin{center}
\begin{tabular}{lr}
\toprule
\textbf{faults occur} & \textbf{alarms that are real} \\
\midrule
1 in 10      & $\val{p23.ppv.10.pct}$\% \\
1 in 100     & \blank \\
1 in 1000    & $\val{p23.ppv.pct}$\% \\
1 in 10\,000 & $\val{p23.ppv.10000.pct}$\% \\
\bottomrule
\end{tabular}
\end{center}
\nextframe
\end{fr}

\begin{fr}
\begin{ansblock}
$\val{p23.ppv.100.pct}$ per cent \dash{} exactly half.
\end{ansblock}

That row is worth carrying away, because it is a landmark rather than a
coincidence. When the fault rate equals the error rate, the real alarms and the
false alarms come from populations whose sizes exactly cancel the two rates,
and the answer is $\frac{1}{2}$ on the nose. It is exact over fractions and the
script asserts it.

So the rule of thumb, and it is the one sentence to keep: \textbf{compare the
base rate with the error rate.} Rarer than the error rate and most alarms are
false; commoner and most are real.
\end{fr}

\begin{fr}
The listing runs the same sweep. It prints fractions rather than decimals,
because the middle entry is exactly one half and a decimal would not say so.

\transcript{p23-base-rate}

The second call is there because a decimal is what you will actually see in a
report, and it is the row that this whole section is about.
\end{fr}

\begin{fr}
One more consequence of the same arithmetic, read from the other end, because
it turns up in evaluation reports constantly.

Our detector is correct on $\val{p23.acc.pct}$ per cent of requests. Consider a
rival that raises no alarm ever, under any circumstances. What is its accuracy?
\dotline
\nextframe
\end{fr}

\begin{fr}
\ans{$\val{p23.never.pct}$ per cent}

It is right about every clean request, and clean requests are
$\val{p23.never.pct}$ per cent of the traffic. \textbf{A model that does
nothing at all scores higher than the working detector}, which is the accuracy
paradox and is the same base rate wearing a different hat.

\begin{warning}
The lesson is not that accuracy is a bad metric. It is that a single number on
an unbalanced population is a weighted average with the weights left out, and
the weights are the base rate. Report the two counts that made it, which is
Program~\ref{prog:F10}'s rule, and the paradox cannot form.
\end{warning}
\end{fr}

% ============================================================
\section{Independence, and what it does not survive}
\label{sec:P23-independence}
\index{probability!independence}

\begin{fr}
Section 4 left the detector looking useless. It is not, and the way out runs
through the next two sections \dash{} but it needs a word this program has been
careful not to use yet.

Two events are \textbf{independent} when
\[ \Prob(A \cap B) = \Prob(A)\,\Prob(B) \]
That is the form to write down, and it is not the form that says what it means.

\begin{notationbox}
\textbf{The word \enquote{independent} is doing two jobs in this book and they
do not meet.}

Program~\ref{prog:P04} spends a program on it: there it is a property of a
\emph{set of vectors}, none of them a combination of the others, and nothing
about it is probabilistic. Here it is a property of \emph{two events under a
measure}, and it is a product of two numbers.

Both are the standard use of the word in their own field, so the collision is
not anybody's mistake and cannot be written around. The defence is to notice
what kind of object the word has been applied to.
\end{notationbox}

Divide both sides by $\Prob(B)$. What does the result say, in words?
\dotline
\nextframe
\end{fr}

\begin{fr}
\ans{$\Prob(A \mid B) = \Prob(A)$: being told $B$ does not move your answer
about $A$}

That is the readable form and the product is still the one to write down,
because the product is symmetric in $A$ and $B$ and does not fall over when
$\Prob(B)$ is zero.

Independence is an assumption about a \emph{measure}, not a fact about two
things being unrelated in the world. Nothing stops you writing it down; whether
it holds is a question about the numbers.

Take two fair coins, so the four outcomes each have weight $\frac{1}{4}$. Let
$A$ be \emph{the first coin is heads} and $B$ \emph{the second is heads}. Are
$A$ and $B$ independent?
\dotline
\nextframe
\end{fr}

\begin{fr}
\ans{Yes}

$\Prob(A) = \Prob(B) = \frac{1}{2}$ and $\Prob(A \cap B) = \frac{1}{4}$, which
is the product. That is what two fair coins ought to do, and it would be a
strange model of coins that said otherwise.

Now add a third event to the same space. Let $C$ be \emph{exactly one of the
two coins is heads}, which is true in two of the four outcomes, so
$\Prob(C) = \frac{1}{2}$.

You are told that $C$ happened. What is $\Prob(A \cap B \mid C)$?
\dotline
\nextframe
\end{fr}

\begin{fr}
\ans{$0$}

If exactly one coin is heads then they cannot both be. Meanwhile
$\Prob(A \mid C)$ and $\Prob(B \mid C)$ are both $\frac{1}{2}$ \dash{} given
$C$, each coin is still equally likely to be the head \dash{} so their product
is $\frac{1}{4}$.

\[ 0 \;\ne\; \tfrac{1}{4} \]

\textbf{$A$ and $B$ are independent, and given $C$ they are as dependent as two
events can be:} knowing either one settles the other completely.
\end{fr}

\begin{fr}
\begin{trapbox}
\textbf{\enquote{These two signals are independent, so I can treat them
separately.}} Independence is not a property two events carry around with them.
It is a statement about a particular measure, and conditioning changes the
measure.

The coins are the smallest possible demonstration and they are not a curiosity.
Every time you filter a dataset, restrict an evaluation to a subgroup, or look
only at the requests that reached the model, you have conditioned \dash{} and
whatever independence held before is now a claim about a different space that
nobody has checked.
\end{trapbox}
\end{fr}

\begin{fr}
The same three events make a smaller point that is worth having separately,
because it comes up whenever more than two things are involved.

Check the other two pairs: $A$ with $C$, and $B$ with $C$. Each pair
multiplies, so all three pairs are independent. What, then, is
$\Prob(A \cap B \cap C)$?
\dotline
\nextframe
\end{fr}

\begin{fr}
\ans{$0$, against a product of $\frac{1}{8}$}

$A \cap B$ is already empty once $C$ is imposed, so adding $C$ cannot rescue
it. \textbf{Pairwise independence does not give you independence of the
collection}, and no amount of checking pairs will find that out.

That is why a claim that a set of features is independent has to say which of
the two it means \dash{} and why a table of pairwise correlations, which is
what people usually produce, cannot settle it.
\end{fr}

% ============================================================
\section{Conditional independence, and a second signal}
\label{sec:P23-conditional}
\index{probability!conditional independence}
\index{probability!likelihood ratio}

\begin{fr}
The last section ran in one direction: independent, and then not, once you
conditioned. It runs the other way just as easily, and that direction is the
one the whole of applied probability is built on.

Two monitors watch the same service. Each is a noisy reading of one underlying
truth \dash{} whether the service is actually broken \dash{} and each is right
three times in four. The two monitors do not talk to each other at all.

Are their two readings independent?
\dotline
\nextframe
\end{fr}

\begin{fr}
\begin{ansblock}
No. $\Prob(X \cap Y) = \frac{\val{p23.cause.pxy.num}}{\val{p23.cause.pxy.den}}$
where $\Prob(X)\Prob(Y) = \frac{1}{4} = \frac{4}{16}$.
\end{ansblock}

They are correlated, and nothing was wired between them. Both are reading the
same underlying state, so when the service is genuinely broken both tend to say
so, and when it is fine both tend to be quiet. A shared cause makes two
independent instruments agree more often than chance.

But now condition on that cause. Given that the service really is broken, the
two readings \emph{are} independent \dash{} exactly, and the script checks it.
Each monitor's noise is its own.
\end{fr}

\begin{fr}
That is \textbf{conditional independence}, and the pair with section 5 is the
thing to carry away:

\begin{center}
\begin{tabular}{ll}
\toprule
\textbf{independent} & and dependent once you condition \\
\textbf{dependent}   & and independent once you condition \\
\bottomrule
\end{tabular}
\end{center}

Both happen, in the smallest examples anybody could build. So neither word
implies the other, and a sentence containing \enquote{independent} that does
not say \emph{given what} has not said anything checkable.

The second row carries a name you have met. Which assumption, made about
several observations and the thing you are trying to infer from them, is it?
\dotline
\nextframe
\end{fr}

\begin{fr}
\ans{The \textbf{naive Bayes} assumption}

Which is more useful stated that way than as the algorithm it is named after,
because it says exactly what is being assumed \dash{} the observations are
conditionally independent given the class \dash{} and therefore exactly what
would break it. Two features that share a cause other than the class break it,
and nothing in the training procedure will report that.

Now cash it in, and undo section 4's damage. Bayes' theorem has a second form
that makes accumulating evidence obvious, and it comes from dividing the
theorem by itself.

Write it for $A$ and again for $\bar{A}$, with the same evidence $B$. The
denominator $\Prob(B)$ is the same in both, so dividing one by the other
cancels it:
\[ \underbrace{\frac{\Prob(A \mid B)}{\Prob(\bar{A} \mid B)}}_{\text{posterior odds}}
   \;=\;
   \underbrace{\frac{\Prob(A)}{\Prob(\bar{A})}}_{\text{prior odds}}
   \times
   \underbrace{\frac{\Prob(B \mid A)}{\Prob(B \mid \bar{A})}}_{\text{likelihood ratio}} \]

\textbf{In odds, the theorem is a multiplication}, and the awkward denominator
has gone.
\end{fr}

\begin{fr}
Put the detector through it. Faults run at one in $\val{p23.prev.den}$, so the
prior odds are $1$ to $999$. An alarm has likelihood ratio
\[ \frac{\Prob(\text{alarm} \mid \text{fault})}
        {\Prob(\text{alarm} \mid \text{clean})} = \val{p23.lr} \]
so one alarm gives posterior odds of $\val{p23.lr}$ to $999$, and a probability
of $\frac{\val{p23.ppv.num}}{\val{p23.ppv.den}}$ \dash{} section 4's answer,
reached by multiplication instead of by a fraction with two terms underneath.

Odds are worth getting comfortable with for exactly that reason. Multiplying a
prior by a ratio is arithmetic you can do while somebody is still talking,
where a fraction with two terms underneath is not; convert to a probability
once, at the end, and only if somebody asks for one.

A second, independent detector of the same quality now raises an alarm as well.
What are the posterior odds, and what probability is that?
\dotline
\nextframe
\end{fr}

\begin{fr}
\begin{ansblock}
Odds of $\val{p23.lr}^{2}$ to $999$, which is
$\frac{\val{p23.two.num}}{\val{p23.two.den}}$, or $\val{p23.two.pct}$ per cent.
\end{ansblock}

$9801$ to $999$ divides by $27$ to give $363$ to $37$, so the probability is
$\frac{363}{400}$. Two alarms take you from $\val{p23.ppv.pct}$ per cent to
$\val{p23.two.pct}$: \textbf{$\val{p23.k.half}$ independent detectors are
enough to cross a half}, and a third leaves only
$\val{p23.three.false.pct}$ per cent of triples false.

That is the constructive answer to section 4. A single weak signal against a
low base rate is nearly worthless, and this is why real monitoring stacks are
built out of several cheap detectors rather than one expensive one.

\transcript{p23-accumulate}
\end{fr}

\begin{fr}
That last step assumed something, and it is worth finding out what by breaking
it. Suppose the second detector is not independent at all but a copy of the
first, keyed on the same log pattern, so that once the first has fired the
second is certain to fire whether or not there is a fault.

What is the second alarm's likelihood ratio, and where does that leave the
answer?
\dotline
\nextframe
\end{fr}

\begin{fr}
\begin{ansblock}
Its ratio is $1$, so it multiplies the odds by nothing and the answer stays at
$\val{p23.ppv.pct}$ per cent.
\end{ansblock}

\begin{trapbox}
\textbf{\enquote{We have two independent signals, so the evidence multiplies.}}
The ratios multiply when the two signals are conditionally independent
\emph{given the truth} \dash{} which is section 5's second row, and which
nothing in either alarm tells you.

The gap between $\val{p23.ppv.pct}$ and $\val{p23.two.pct}$ per cent
\textbf{is} what the independence assumption is worth, and it is the whole of
it. Two detectors sharing a feature, two evaluations sharing prompts, two
reviewers who read each other's comments: each is the duplicate case wearing
different clothes, and each will be reported as corroboration.
\end{trapbox}
\end{fr}

\begin{fr}
\begin{rigourbox}
This program treats a probability as a weight on a finite list of outcomes,
which is enough for everything in it and is not the general definition. A
proper treatment names the collection of measurable sets and requires
additivity over countably many disjoint pieces, which is what makes continuous
spaces work \dash{} where, as Program~\ref{prog:F13} says, a quantity with
infinitely many possible values cannot carry a probability at each one, and
areas do the work instead. A first course in measure theory
does it properly; nothing in this book turns on the difference.

What this program also does not do is say what a probability \emph{means}.
The three rules are silent on it, and the two readings people hold \dash{} a
long-run frequency, or a degree of belief \dash{} agree on every calculation in
this book and disagree about what to do when there is no long run. That
argument is not settled here. Program~\ref{prog:P27} states what a p-value does
and does not say, and Program~\ref{prog:P28} contrasts a credible interval with
a confidence interval, which is where the two readings answer the same question
differently.
\end{rigourbox}
\end{fr}

\begin{fr}
Where this leaves you. You can build a sample space, put a measure on it and
check that it is one; you can condition, which is choosing a denominator on
purpose; you can derive Bayes' theorem rather than remembering it; and you can
say which of the two conditionals a quoted accuracy is, which is the single
most valuable thing in the program.

What you do not have yet is a way to summarise a distribution rather than
compute with it. Program~\ref{prog:P24} makes a random variable a function on
this sample space and gives it two summaries \dash{} which is also where the
sampling a language model actually does turns out to be a draw from a
distribution you have been editing all along.
\end{fr}

\begin{summarybox}
\sumitem{1--3}{A \textbf{sample space} is the list of outcomes, exactly one of
  which happens; an \textbf{event} is a set of them. A probability hands each
  event a number that is never negative, gives the whole space $1$, and adds
  over pieces that cannot overlap.}
\sumitem{4--5}{$\Prob(\bar{A}) = 1 - \Prob(A)$, from two of the three rules and
  no picture. Every result here is got the same way: split into disjoint
  pieces, add, use that the whole is $1$.}
\sumitem{6--7}{Additivity needs \emph{disjoint} pieces:
  $\Prob(A \cup B) = \Prob(A) + \Prob(B) - \Prob(A \cap B)$. That is
  Program~\ref{prog:F10}'s union rule with every count divided by
  $\val{f10.eval.union}$, and the script gates the two against each other.}
\sumitem{8--9}{A probability is a \textbf{measure}, and rule two bounds the
  measure of a \emph{set}. Program~\ref{prog:F13}'s density height is a rate
  rather than the measure of a set, which is why it may exceed $1$ where a
  probability may not.}
\sumitem{10}{\textbf{Conditioning is changing the denominator}:
  $\Prob(A \mid B) = \Prob(A \cap B)/\Prob(B)$. Two of the four outcomes are
  simply gone, and what is left is renormalised.}
\sumitem{11--12}{It is a \textbf{definition} and not a theorem, which is
  Program~\ref{prog:F10}'s decision made on purpose. Bayes' theorem is the
  other kind, and that is what makes it rebuildable.}
\sumitem{13}{Of $\val{p23.alarm.den.long}$ requests,
  $\val{p23.alarm.num}/\val{p23.alarm.den}$ produce an alarm \dash{}
  $\val{p23.alarm.real}$ real and $\val{p23.alarm.false}$ false. Keep the two
  counts, not the total.}
\sumitem{14--15}{$\Prob(A \cap B) = \Prob(A \mid B)\Prob(B)$ is the definition
  with the denominator moved, $\Prob(B) > 0$ is the only condition anywhere in
  the program, and $\Prob(A \mid B)$ and $\Prob(B \mid A)$ are different
  numbers.}
\sumitem{16--17}{\textbf{Bayes' theorem is one quantity written two ways and
  then divided.} Nothing is assumed and no new object appears; it reverses a
  conditional, which is the shape of most measurement problems here.}
\sumitem{18--19}{Prior, likelihood, evidence, posterior \dash{} names that will
  not help you rebuild the formula. The derivation will.}
\sumitem{20--21}{$\Prob(B) = \Prob(B \mid A)\Prob(A) +
  \Prob(B \mid \bar{A})\Prob(\bar{A})$, the law of total probability. It
  contains $\Prob(A)$, so the base rate is in the answer whether or not
  anybody mentioned it.}
\sumitem{22--23}{A detector $\val{p23.sens.pct}$ per cent accurate on a fault
  once in $\val{p23.prev.den}$ requests has
  $\frac{\val{p23.ppv.num}}{\val{p23.ppv.den}}$ of its alarms real \dash{}
  $\val{p23.ppv.pct}$ per cent, about $\val{p23.alarm.ratio}$ false alarms for
  every real one.}
\sumitem{23--24}{$\val{p23.sens.pct}$ per cent is
  $\Prob(\text{alarm} \mid \text{fault})$ and the question was
  $\Prob(\text{fault} \mid \text{alarm})$. A detector's quoted numbers are
  properties of the detector \emph{and} the population, and only one of the
  two travels with the model.}
\sumitem{25--26}{The \textbf{prosecutor's fallacy} is swapping those two:
  $\val{p23.fpr.pct}$ per cent of clean requests raise an alarm and
  $\val{p23.false.pct}$ per cent of alarms are about clean requests, both true
  of the same detector on the same day.}
\sumitem{26--28}{Same detector, four base rates:
  $\val{p23.ppv.10.pct}$\%, $\val{p23.ppv.100.pct}$\%, $\val{p23.ppv.pct}$\%,
  $\val{p23.ppv.10000.pct}$\%. \textbf{Compare the base rate with the error
  rate}: equal gives exactly half, and the transcript prints it as a fraction
  so that it says so.}
\sumitem{29--30}{A detector that never fires is $\val{p23.never.pct}$ per cent
  accurate against the real one's $\val{p23.acc.pct}$. Accuracy on an
  unbalanced population is a weighted average with the weights left out.}
\sumitem{31--32}{$A$ and $B$ are independent when
  $\Prob(A \cap B) = \Prob(A)\Prob(B)$, which divided through says being told
  $B$ does not move your answer about $A$.}
\sumitem{33--35}{Two fair coins are independent and, given that exactly one is
  heads, perfectly dependent \dash{} \textbf{independence is a statement about
  a measure, and conditioning changes the measure}. Every filter, subgroup and
  \enquote{requests that reached the model} is a conditioning.}
\sumitem{36--37}{Pairwise independence is not independence of the collection:
  those same three events are independent in every pair and have
  $\Prob(A \cap B \cap C) = 0$ against a product of $\frac{1}{8}$.}
\sumitem{38--39}{Two noisy readings of one cause are \emph{dependent}
  ($\frac{\val{p23.cause.pxy.num}}{\val{p23.cause.pxy.den}}$ against
  $\frac{4}{16}$) with nothing wired between them, and conditionally
  independent given the cause.}
\sumitem{40--41}{Both directions happen, so neither word implies the other. The
  second is the \textbf{naive Bayes} assumption, and in odds Bayes' theorem is
  a multiplication: posterior odds are prior odds times the likelihood ratio.}
\sumitem{42--43}{One alarm at ratio $\val{p23.lr}$ gives
  $\val{p23.ppv.pct}$ per cent; two give
  $\frac{\val{p23.two.num}}{\val{p23.two.den}} = \val{p23.two.pct}$ per cent,
  so $\val{p23.k.half}$ independent detectors cross a half and a third leaves
  $\val{p23.three.false.pct}$ per cent of triples false.}
\sumitem{44--45}{A duplicate detector has likelihood ratio $1$ and moves
  nothing, so the gap between $\val{p23.ppv.pct}$ and $\val{p23.two.pct}$ per
  cent is exactly what the independence assumption is worth.}
\sumitem{46--47}{Not proved here: countable additivity, and what a probability
  \emph{means}. Neither changes a calculation in this book.}
\end{summarybox}

\canyou

\begin{testexercises}
\item A screening test is right $\val{p23.sens.pct}$ per cent of the time in
  both directions, and the condition it screens for affects one person in
  $\val{p23.prev.den}$. Somebody tests positive. Without computing, say
  whether the probability they have the condition is nearer
  $\val{p23.sens.pct}$ per cent or nearer $\val{p23.ppv.pct}$ per cent, and
  say why in one sentence.
  \answerto{Nearer $\val{p23.ppv.pct}$ per cent. The condition is far rarer
  than the test's error rate, so the false positives drawn from the enormous
  healthy population outnumber the true positives drawn from the tiny affected
  one. The quoted figure is
  $\Prob(\text{positive} \mid \text{condition})$; the question asks for
  $\Prob(\text{condition} \mid \text{positive})$.}
\item Two events satisfy $\Prob(A) = \Prob(B) = \frac{1}{2}$ and
  $\Prob(A \cap B) = \frac{1}{3}$. Are they independent, and what is
  $\Prob(A \cup B)$?
  \answerto{Not independent: the product is $\frac{1}{4}$ and the joint is
  $\frac{1}{3}$, so knowing one raises the chance of the other. The union is
  $\frac{1}{2} + \frac{1}{2} - \frac{1}{3} = \frac{2}{3}$.}
\item Your team reports that a classifier has precision $\val{p23.sens.pct}$
  per cent on the benchmark and is unusable in production. Name the quantity
  that differs between the two settings, and say what you would ask for to
  confirm it.
  \answerto{The base rate. A benchmark is usually balanced or close to it and
  production traffic is not, and precision is
  $\Prob(\text{true} \mid \text{flagged})$, which depends on the population as
  well as on the model. Ask for the two counts behind each figure \dash{} how
  many cases were flagged and how many of those were real \dash{} in the
  benchmark and in production separately.}
\item Prior odds on a hypothesis are $1$ to $\val{p23.prev.den}$. Two
  independent tests each have likelihood ratio $\val{p23.lr}$ and both come
  back positive. What are the posterior odds, and roughly what probability is
  that?
  \answerto{Multiply: $\val{p23.lr}^{2} = 9801$ against $\val{p23.prev.den}$,
  so odds of about $9801$ to $\val{p23.prev.den}$, which is a probability just
  under $\frac{10}{11}$ \dash{} around ninety per cent. The prior odds here
  are $1$ to $\val{p23.prev.den}$ rather than $1$ to $999$, so the answer is
  slightly below the $\val{p23.two.pct}$ per cent in the text.}
\item Give an example of two events that are independent and become dependent
  when a third is given, and one that runs the other way. One sentence each.
  \answerto{Two fair coins are independent, and given that exactly one landed
  heads each determines the other. Two noisy readings of one underlying state
  are dependent, and given that state they are independent, because each
  reading's noise is its own.}
\end{testexercises}

\begin{furtherproblems}
\item Show that if $A$ and $B$ are independent then so are $A$ and $\bar{B}$.
  \answerto{$\Prob(A \cap \bar{B}) = \Prob(A) - \Prob(A \cap B)
  = \Prob(A) - \Prob(A)\Prob(B) = \Prob(A)(1 - \Prob(B))
  = \Prob(A)\Prob(\bar{B})$. The first step splits $A$ by whether $B$ happened,
  which is two disjoint pieces, so rule three applies.}
\item At what fault rate does a detector with sensitivity $s$ and false-positive
  rate $f$ have exactly half its alarms real? Check your answer against the
  program's numbers.
  \answerto{Real alarms are $ps$ and false ones $(1-p)f$; setting them equal
  gives $p = f/(s+f)$. With $s = \num{0.99}$ and $f = \num{0.01}$ that is
  $\num{0.01}$, one in a hundred, which is the row the table asks you to fill
  in. Note the answer does not depend on $s$ and $f$ separately, only on their
  ratio.}
\item A detector's alarm has likelihood ratio $L$. How many independent alarms
  are needed before the posterior odds exceed $1$, starting from prior odds
  $1$ to $N$?
  \answerto{You need $L^{k} > N$, so $k > \logb{10} N / \logb{10} L$ and the
  answer is the next integer above it. With $N = 999$ and
  $L = \val{p23.lr}$ that is $k > \num{1.503}$, so
  $k = \val{p23.k.half}$ \dash{} which is the number the text measures. The
  base rate enters logarithmically, so a hundredfold rarer fault needs only
  one more independent detector.}
\item Three events are pairwise independent and satisfy
  $\Prob(A \cap B \cap C) = 0$. Construct such a triple other than the coins in
  section 5, or argue that the coins are essentially the only one on four
  outcomes.
  \answerto{On four equally weighted outcomes the coins are the construction:
  take any three events of probability $\frac{1}{2}$ whose pairwise
  intersections each contain exactly one outcome and which have no common
  outcome. Relabelling the outcomes gives every such triple, so up to renaming
  it is the only one. On more outcomes there are others \dash{} three of the
  six faces of a die, for instance, chosen so no face lies in all three.}
\item Section 6 says a duplicate detector has likelihood ratio $1$. Write down
  what would have to be true of the two detectors for the ratio of the second
  to be exactly $1$, and say whether that is the same as the two detectors
  being identical.
  \answerto{The second alarm must be equally likely whether or not there is a
  fault, \emph{given the first alarm} \dash{} that is,
  $\Prob(A_2 \mid A_1, F) = \Prob(A_2 \mid A_1, \bar{F})$. Identical detectors
  are one way to arrange it and not the only one: any second detector whose
  extra information is entirely contained in the first has this property, and
  two detectors can look quite different and still share the one feature they
  are both really keyed on.}
\end{furtherproblems}
